% !TeX program = xelatex
% Informe de Perfil de Proyecto - Sistema IoT para terma doméstica
\documentclass[11pt,a4paper]{article}

\usepackage[utf8]{inputenc}
\usepackage[T1]{fontenc}
\usepackage[spanish,es-nodecimaldot]{babel}
\usepackage{geometry}
\geometry{a4paper, left=2.4cm, right=2.4cm, top=2.6cm, bottom=2.4cm}
\usepackage{graphicx}
\usepackage{amsmath,amssymb}
\usepackage{booktabs}
\usepackage{tabularx}
\usepackage{array}
\usepackage{ragged2e}
\newcolumntype{L}{>{\RaggedRight\arraybackslash}X}
\newcolumntype{P}[1]{>{\RaggedRight\arraybackslash}p{#1}}
\usepackage{enumitem}
\usepackage{xcolor}
\usepackage{titlesec}
\usepackage{fancyhdr}
\usepackage{caption}
\usepackage{tikz}
\usetikzlibrary{arrows.meta, positioning, babel}
\usepackage[hidelinks]{hyperref}

% ---------- Metadatos ----------
\hypersetup{
  pdftitle={Perfil de Proyecto: Sistema IoT Inteligente para el Monitoreo, Control Remoto y Gestion Energetica de una Terma Domestica mediante Asistente Virtual},
  pdfauthor={Proyecto Terma IoT},
  pdfsubject={Perfil de Proyecto - Ingenieria Electronica / IoT},
  pdfkeywords={IoT, ESP32, terma, eficiencia energetica, MQTT, asistente virtual}
}

% ---------- Color ----------
\definecolor{primary}{HTML}{0E4C5C}
\definecolor{accent}{HTML}{C77700}
\definecolor{softline}{HTML}{B9C6CC}
\definecolor{codebg}{HTML}{F2F5F6}

% ---------- Estilo de secciones ----------
\titleformat{\section}{\normalfont\Large\bfseries\color{primary}}{\thesection.}{0.6em}{}
\titleformat{\subsection}{\normalfont\large\bfseries\color{primary!85!black}}{\thesubsection}{0.6em}{}
\titlespacing*{\section}{0pt}{1.4em}{0.6em}
\titlespacing*{\subsection}{0pt}{1.0em}{0.4em}

\captionsetup{font=small,labelfont={bf,color=primary},justification=centering}
\setlist[itemize]{leftmargin=1.4em,itemsep=0.25em,topsep=0.3em}
\setlist[enumerate]{leftmargin=1.6em,itemsep=0.25em,topsep=0.3em}

\pagestyle{fancy}
\fancyhf{}
\fancyhead[L]{\footnotesize\color{primary}Perfil de Proyecto --- Terma Doméstica Inteligente (IoT)}
\fancyhead[R]{\footnotesize\thepage}
\renewcommand{\headrulewidth}{0.4pt}
\renewcommand{\headrule}{\hbox to\headwidth{\color{softline}\leaders\hrule height \headrulewidth\hfill}}

\newcommand{\tblhead}[1]{\textbf{\color{primary}#1}}

\begin{document}

% ============================================================
% PORTADA
% ============================================================
\begin{titlepage}
\centering
\vspace*{1.2cm}
{\color{softline}\rule{\textwidth}{0.8pt}}\\[0.5cm]
{\large\scshape Perfil de Proyecto}\\[0.35cm]
{\color{accent}\rule{0.28\textwidth}{2.2pt}}\\[1.0cm]

{\Huge\bfseries\color{primary} Sistema IoT Inteligente para el Monitoreo, Control Remoto y Gestión Energética de una Terma Doméstica mediante Asistente Virtual\par}

\vspace{1.2cm}
{\large Curso: \rule{6cm}{0.4pt}}\\[0.45cm]
{\large Docente: \rule{6cm}{0.4pt}}\\[0.45cm]
{\large Integrantes: \rule{6cm}{0.4pt}}\\[0.25cm]
{\large \phantom{Integrantes:} \rule{6cm}{0.4pt}}\\[0.45cm]
{\large Institución: \rule{6cm}{0.4pt}}\\[0.45cm]
{\large Fecha: \rule{6cm}{0.4pt}}\\[1.4cm]

{\color{softline}\rule{\textwidth}{0.8pt}}

\vfill
\begin{minipage}{0.86\textwidth}
\small
\textbf{\color{primary}Resumen.} Se presenta el perfil de un proyecto cuyo objetivo es diseñar e implementar un prototipo de sistema de Internet de las Cosas (IoT) para el monitoreo, control remoto y gestión energética de una terma eléctrica doméstica. El sistema integra tres sensores (temperatura del agua, corriente eléctrica no invasiva y flujo de agua), un actuador de potencia de estado sólido (SSR), comunicación inalámbrica mediante un microcontrolador ESP32, un dashboard con procesamiento previo de datos y una interfaz opcional de asistente virtual por voz. El documento desarrolla la introducción, la problemática, el marco teórico, el estado del arte (con ocho artículos científicos), la metodología (componentes, microcontrolador y diagrama de bloques) y los objetivos y alcances del prototipo.
\end{minipage}
\end{titlepage}

\pagenumbering{roman}
\tableofcontents
\clearpage
\pagenumbering{arabic}

% ============================================================
\section{Introducción}
% ============================================================
El Internet de las Cosas (IoT) ha transformado la manera en que los dispositivos cotidianos miden, comunican y actúan sobre su entorno. En el ámbito residencial, esta tecnología habilita los denominados \emph{hogares inteligentes}, capaces de supervisar variables físicas, automatizar decisiones y entregar información útil al usuario en tiempo real~\cite{alfuqaha2015}. Dentro de este ecosistema, la gestión energética doméstica (\emph{Home Energy Management}, HEM) constituye una de las líneas de mayor impacto, pues permite reducir el consumo, evitar desperdicios y mejorar el confort del habitante~\cite{ford2017,zipperer2013}.

En el contexto local, la terma eléctrica (calentador de agua por resistencia) es uno de los artefactos de mayor demanda energética del hogar, dado que transforma energía eléctrica en calor mediante resistencias de potencias típicas del orden de 1{,}5 a 3 kW. A diferencia de otros equipos, su operación suele depender de un interruptor manual y de la memoria del usuario: la terma se enciende ``por si acaso'' o permanece activa sin que exista certeza de que el agua vaya a utilizarse. Esa forma de operación conduce a un consumo innecesario y a la ausencia de información objetiva sobre el estado del equipo.

Este proyecto propone desarrollar un prototipo funcional de \textbf{terma doméstica inteligente} basado en un microcontrolador con conectividad inalámbrica (ESP32), capaz de: (i) medir la temperatura del agua, la corriente eléctrica consumida y el flujo de agua; (ii) accionar remotamente la carga mediante un relé de estado sólido; (iii) transmitir los datos a un dashboard que calcula métricas energéticas; (iv) ejecutar reglas automáticas de encendido y apagado; y (v) ofrecer, de forma opcional, comandos por voz a través de un asistente virtual.

El documento se organiza de la siguiente manera. La Sección~\ref{sec:problematica} describe y justifica la problemática. La Sección~\ref{sec:marco} expone el marco teórico. La Sección~\ref{sec:arte} analiza el estado del arte a partir de ocho artículos científicos. La Sección~\ref{sec:metodologia} detalla la metodología, la lista de componentes, el microcontrolador seleccionado y el diagrama de bloques. Finalmente, la Sección~\ref{sec:objetivos} precisa los objetivos y los alcances y limitaciones del prototipo.

% ============================================================
\section{Problemática}
\label{sec:problematica}
% ============================================================

\subsection*{Descripción del problema}
Las termas eléctricas domésticas operan, en la mayoría de los casos, bajo un control estrictamente manual: el usuario enciende el equipo y, posteriormente, debe recordar apagarlo. Esta dependencia humana genera dos fallas recurrentes: (i) dejar la terma encendida durante periodos prolongados, incluso cuando el agua no se utiliza, y (ii) encenderla con una anticipación excesiva respecto del momento real de uso. En ambos escenarios la resistencia consume energía eléctrica sin un beneficio directo para el usuario, elevando la facturación y el impacto ambiental asociado~\cite{bhati2017}.

A lo anterior se suma la falta de retroalimentación. El usuario no dispone de información en tiempo real sobre la temperatura del agua, la potencia que está consumiendo el equipo ni si éste se encuentra efectivamente en uso. Sin datos, resulta imposible tomar decisiones informadas de ahorro: el comportamiento del usuario frente a la energía suele mejorar cuando dispone de mediciones visibles y comprensibles~\cite{zipperer2013,ford2017}. Diversos estudios muestran, además, que la adopción de tecnología en el hogar depende de que la solución sea percibida como útil, confiable y poco intrusiva para el usuario final~\cite{marikyan2019}.

Finalmente, la propia configuración eléctrica de la terma introduce requisitos de seguridad que el control manual no atiende: el accionamiento de una carga resistiva de alta potencia, sumergida y en contacto con agua, exige aislamiento galvánico y una conmutación robusta. Un sistema IoT que controle esta carga debe garantizar la separación entre la electrónica de baja tensión y la red de 220 V AC.

\subsection*{Justificación}
El proyecto se justifica en tres dimensiones complementarias:

\begin{itemize}
  \item \textbf{Económica:} reducir el consumo innecesario de una de las cargas de mayor potencia del hogar, aprovechando el monitoreo para encender sólo cuando se requiere y apagar de forma automática al alcanzar la temperatura objetivo o tras un periodo de inactividad~\cite{bhati2017}.
  \item \textbf{Tecnológica:} demostrar que es posible construir una solución funcional y de bajo costo con componentes estándar (ESP32, sensores DS18B20, SCT-013 y YF-S201, y un SSR), sin recurrir a procesadores costosos.
  \item \textbf{Social y de sostenibilidad:} aportar información objetiva al usuario para promover hábitos de ahorro energético y contribuir a la reducción del consumo residencial, alineado con la literatura de gestión energética del hogar~\cite{ford2017,zipperer2013,marikyan2019}.
\end{itemize}

En consecuencia, la problemática puede resumirse así: \emph{la operación manual y no supervisada de la terma doméstica produce un consumo energético innecesario y una experiencia de usuario deficiente, situación que un sistema IoT de bajo costo ---con sensado múltiple, actuación segura, comunicación inalámbrica y visualización con procesamiento previo--- puede corregir de forma efectiva.}

% ============================================================
\section{Marco Teórico}
\label{sec:marco}
% ============================================================

\subsection{Internet de las Cosas y arquitecturas de referencia}
El IoT se define como la interconexión de objetos físicos dotados de capacidad de sensado, cómputo y comunicación, que intercambian datos a través de redes, frecuentemente inalámbricas, para ofrecer servicios a aplicaciones y usuarios~\cite{alfuqaha2015}. Una arquitectura de referencia organiza el sistema en capas: \textbf{percepción} (sensores y actuadores), \textbf{red} (transporte de datos) y \textbf{aplicación} (almacenamiento, análisis y presentación). El proyecto adopta este modelo de tres capas, con un componente de \emph{procesamiento previo} (\emph{edge preprocessing}) ubicado en el propio nodo, que filtra, calibra y agrega los datos antes de transmitirlos. Esta estrategia reduce el volumen de información enviada y mejora la latencia y la resiliencia de la aplicación~\cite{shi2016}.

\subsection{Protocolos de comunicación}
Para el transporte de datos sobre Wi-Fi, el protocolo \emph{Message Queuing Telemetry Transport} (MQTT) es ampliamente utilizado en aplicaciones IoT por su modelo publicador/suscriptor, su bajo consumo de ancho de banda y su tolerancia a enlaces inestables. Frente a alternativas basadas en petición/respuesta (HTTP), MQTT resulta más eficiente para telemetría periódica de muchos dispositivos~\cite{alfuqaha2015}. El sistema propuesto publica las lecturas en tópicos MQTT y recibe comandos de control en tópicos de actuación.

\subsection{Medición de corriente y desagregación de cargas}
La medición de corriente eléctrica puede realizarse de forma invasiva (con derivación o \emph{shunt}) o no invasiva, mediante transformadores de corriente (CT) tipo pinza que se instalan alrededor del conductor sin interrumpir el circuito. Los CT permiten medir corriente alterna de manera segura y económica. La señal alterna obtenida se procesa típicamente calculando su valor eficaz (\emph{RMS}) y, a partir de éste, la potencia aparente o activa. El estudio de la desagregación de cargas no intrusiva (NILM) muestra que la firma eléctrica de un equipo puede caracterizarse a partir de su corriente y potencia para identificar su funcionamiento~\cite{zoha2012}. En este proyecto, la firma de la terma (carga resistiva prácticamente constante mientras calienta) permite inferir su estado de operación a partir de la corriente medida.

Las magnitudes energéticas básicas empleadas son la potencia eléctrica
\begin{equation}
P(t) = V_{\mathrm{rms}}(t)\,\cdot\, I_{\mathrm{rms}}(t)\,\cdot\, \cos\varphi ,
\end{equation}
y la energía consumida en un intervalo, obtenida por integración:
\begin{equation}
E = \int_{t_0}^{t_1} P(t)\,\mathrm{d}t \;\;\approx\;\; \sum_{k} P_k \,\Delta t .
\end{equation}
Para una carga resistiva pura, el factor de potencia $\cos\varphi \approx 1$, lo que simplifica el cálculo de la energía a partir de la corriente y la tensión de red.

\subsection{Sensado de temperatura}
El termómetro digital DS18B20 es un sensor de temperatura con interfaz digital 1-Wire que entrega la medición en formato digital, con un rango de operación de $-55\,^{\circ}\mathrm{C}$ a $+125\,^{\circ}\mathrm{C}$ y una exactitud de $\pm 0{,}5\,^{\circ}\mathrm{C}$ en el intervalo de $-10\,^{\circ}\mathrm{C}$ a $+85\,^{\circ}\mathrm{C}$, con resolución configurable de 9 a 12 bits~\cite{ds18b20}. Su versión con sonda sumergible y empaque de acero inoxidable es adecuada para medir la temperatura del agua de la terma. La comunicación digital 1-Wire permite conectar varios sensores sobre un mismo bus, lo que facilita la redundancia de medición.

\subsection{Medición de flujo de agua}
El caudal puede medirse con sensores de efecto Hall de tipo turbina, como el YF-S201, que genera una señal de pulsos cuya frecuencia es proporcional al flujo volumétrico (típicamente alrededor de 450 pulsos por litro, con un rango de 1 a 30 L/min)~\cite{yfs201201}. Al contar los pulsos por unidad de tiempo se obtiene el caudal instantáneo, y al integrarlo, el volumen consumido. Esta variable permite distinguir entre ``terma encendida sin uso'' y ``terma en uso'', habilitando reglas automáticas basadas en la demanda real de agua.

\subsection{Actuación de potencia}
El control de una carga resistiva de alta potencia se realiza mediante un relé de estado sólido (SSR), dispositivo basado en semiconductores (típicamente optoacoplados) que conmuta la carga AC con aislamiento galvánico entre el circuito de control de baja tensión y la red. Frente a los relés electromecánicos, el SSR ofrece mayor vida útil, ausencia de rebote mecánico y conmutación por cruce por cero (\emph{zero-cross}), lo que reduce el ruido electromagnético~\cite{ssr}. El SSR constituye el actuador principal del sistema y es el elemento que permite el encendido y apagado remoto y automático de la terma.

\subsection{Procesamiento previo y cómputo en el borde}
El \emph{edge computing} propone procesar los datos cerca de su origen para disminuir el tráfico, la latencia y la dependencia de la nube~\cite{shi2016}. En el proyecto, el procesamiento previo incluye: filtrado digital de la señal de corriente (promediado y cálculo de RMS), calibración de sensores, suavizado de la temperatura (media móvil), conteo e integración de pulsos de flujo, y evaluación local de reglas de seguridad (por ejemplo, corte de potencia ante sobretemperatura). Sólo las métricas agregadas se publican en la nube.

% ============================================================
\section{Estado del Arte}
\label{sec:arte}
% ============================================================
En esta sección se analizan ocho trabajos científicos directamente vinculados al proyecto. Los primeros cuatro abordan la gestión energética del hogar y la adopción de tecnología por parte del usuario; los siguientes, la medición y desagregación de cargas eléctricas y las implementaciones de bajo costo basadas en ESP32.

\subsection*{A1. Revisión de aplicaciones de hogar inteligente basadas en IoT}
Alaa y colaboradores~\cite{alaa2017} realizan una revisión sistemática del panorama del hogar inteligente basado en IoT, construyendo una taxonomía de aplicaciones y señalando que la investigación se encuentra dispersa y con desafíos abiertos en interoperabilidad, costo y experiencia de usuario. El trabajo fundamenta la pertinencia de soluciones integradas y de bajo costo como la propuesta, y aporta un marco para clasificar la aplicación (monitoreo + control + gestión).

\subsection*{A2. Revisión de IoT para el hogar inteligente: retos y soluciones}
Stojkoska y Trivodaliev~\cite{stojkoska2017} revisan las tecnologías habilitantes del hogar inteligente y sus retos ---conectividad, seguridad, consumo energético de los nodos y heterogeneidad de dispositivos---, proponiendo que la arquitectura y la selección de protocolos son decisiones críticas. De este trabajo se desprende la conveniencia de usar un microcontrolador con Wi-Fi integrado y un protocolo ligero como MQTT.

\subsection*{A3. Conservación de energía mediante hogares inteligentes}
Bhati, Hansen y Chan~\cite{bhati2017} estudian el impacto de las tecnologías de hogar inteligente en la conservación de energía en hogares de Singapur, mostrando que la retroalimentación (\emph{feedback}) y la automatización contribuyen al ahorro, y que el efecto se potencia con la integración entre dispositivos, hogar y red eléctrica. Este antecedente sustenta la hipótesis de ahorro del proyecto y la importancia de visualizar la información al usuario.

\subsection*{A4. Categorías y funcionalidad de la tecnología de hogar inteligente para la gestión energética}
Ford, Pritoni, Sanguinetti y Karlin~\cite{ford2017} analizan 308 productos comerciales de gestión energética del hogar y clasifican su funcionalidad, evidenciando la brecha entre las capacidades anunciadas y las realmente útiles. El trabajo justifica priorizar un conjunto acotado de funciones bien implementadas (monitoreo, control remoto, reglas automáticas) por encima de la cantidad de prestaciones.

\subsection*{A5. Revisión sistemática de la literatura de hogar inteligente: la perspectiva del usuario}
Marikyan, Papagiannidis y Alamanos~\cite{marikyan2019} revisan la literatura desde la perspectiva del usuario, identificando factores como utilidad percibida, confianza, privacidad y facilidad de uso como determinantes de la adopción. Este antecedente orienta las decisiones de diseño de la interfaz y del asistente virtual, que deben ofrecer interacción simple y comprensible.

\subsection*{A6. Enfoques de monitoreo no intrusivo de cargas}
Zoha, Gluhak, Imran y Rajasegarar~\cite{zoha2012} presentan una revisión de las técnicas de monitoreo no intrusivo de cargas (NILM), describiendo cómo la corriente y la potencia permiten identificar el estado y la operación de los equipos. El proyecto aplica esta idea a una carga única y bien caracterizada (la terma), usando su firma electrica para inferir calentamiento y detectar anomalías.

\subsection*{A7. Gestión de energía eléctrica en el hogar inteligente}
Zipperer y colaboradores~\cite{zipperer2013} analizan las tecnologías habilitantes y el comportamiento del consumidor en la gestión energética del hogar, concluyendo que el éxito depende tanto del componente técnico como de la disposición del usuario a interactuar con la información. Refuerza la necesidad de soluciones que combinen medición, control y una interfaz amigable.

\subsection*{A8. Sistema de monitoreo de energía inteligente basado en ESP32}
El-Khozondar y colaboradores~\cite{elkhozondar2024} desarrollan un sistema de monitoreo energético de bajo costo basado en el microcontrolador ESP32, integrando sensores de tensión (ZMPT101B) y de corriente (SCT-013), y reportando los datos al usuario a través de la red, demostrando la viabilidad del ESP32 y de los transformadores de corriente para el monitoreo eléctrico doméstico. Este trabajo es el antecedente más cercano al componente de medición eléctrica del proyecto.

\subsection*{Síntesis comparativa y vacío identificado}
La Tabla~\ref{tab:arte} resume los trabajos analizados. Se observa que la literatura aborda, por un lado, la gestión energética y la adopción del usuario (A1--A5) y, por otro, la medición de cargas y las implementaciones de bajo costo (A6--A8). Sin embargo, existe un vacío en soluciones \emph{específicas para la terma doméstica} que integren simultáneamente: tres variables físicas (temperatura, corriente y flujo), un actuador de potencia seguro (SSR), comunicación inalámbrica con procesamiento previo y una interfaz que combine dashboard y asistente virtual. El presente proyecto se posiciona en ese vacío.

\begin{table}[htbp]
\centering
\caption{Síntesis comparativa del estado del arte.}
\label{tab:arte}
\footnotesize
\begin{tabularx}{\textwidth}{@{}P{0.7cm} P{3.1cm} L P{2.5cm}@{}}
\toprule
\tblhead{\#} & \tblhead{Trabajo} & \tblhead{Aporte principal} & \tblhead{Limitación / brecha} \\
\midrule
A1 & Alaa et al. (2017) & Taxonomía de aplicaciones del hogar IoT & No enfocado en una carga específica \\
A2 & Stojkoska y Trivodaliev (2017) & Retos de conectividad y protocolos & Sin implementación de control de carga \\
A3 & Bhati et al. (2017) & Evidencia de ahorro por feedback y automatización & Dependiente del contexto del hogar \\
A4 & Ford et al. (2017) & Categorías funcionales de HEM (308 productos) & Sin desarrollo de prototipo \\
A5 & Marikyan et al. (2019) & Factores de adopción desde el usuario & Revisión, no sistema técnico \\
A6 & Zoha et al. (2012) & Métodos NILM para identificar cargas & Alta complejidad de cómputo \\
A7 & Zipperer et al. (2013) & Tecnologías habilitantes y comportamiento & Sin prototipo de bajo costo \\
A8 & El-Khozondar et al. (2024) & Monitoreo eléctrico con ESP32 y SCT-013 & Sólo monitoreo eléctrico \\
\bottomrule
\end{tabularx}
\end{table}

% ============================================================
\section{Metodología}
\label{sec:metodologia}
% ============================================================

\subsection{Arquitectura general del sistema}
El sistema sigue una arquitectura de tres capas con procesamiento previo en el nodo. En la \textbf{capa de percepción}, tres sensores adquieren las variables físicas. En la \textbf{capa de red/proceso}, el ESP32 acondiciona, filtra y procesa las señales, evalúa reglas locales y se comunica por Wi-Fi mediante MQTT. En la \textbf{capa de aplicación}, un dashboard centraliza la visualización y el cálculo de métricas, y un asistente virtual permite el control por lenguaje natural. El SSR actúa sobre la carga (la terma) de forma segura y aislada.

\subsection{Lista de componentes y especificaciones técnicas}
La Tabla~\ref{tab:componentes} presenta la lista de componentes. Se cumple con el mínimo exigido: \textbf{tres sensores} (temperatura, corriente y flujo), \textbf{un actuador} de potencia (SSR), \textbf{comunicación inalámbrica} (Wi-Fi del ESP32), \textbf{interfaz con procesamiento previo} (dashboard) y \textbf{enclosure} de protección.

\begin{table}[htbp]
\centering
\caption{Componentes y especificaciones técnicas del prototipo.}
\label{tab:componentes}
\footnotesize
\begin{tabularx}{\textwidth}{@{}P{3.0cm} P{2.8cm} L@{}}
\toprule
\tblhead{Componente} & \tblhead{Modelo} & \tblhead{Especificación técnica clave y función} \\
\midrule
Microcontrolador & ESP32-WROOM-32 & Dual-core Xtensa LX6 a 240 MHz, 520 KB SRAM, 4 MB flash, Wi-Fi 802.11 b/g/n (2,4 GHz) y BLE 4.2, ADC de 12 bits. Adquisición, procesamiento previo y comunicación. \\
Sensor de temperatura & DS18B20 (sonda sumergible) & Rango $-55$ a $+125\,^{\circ}\mathrm{C}$, exactitud $\pm 0{,}5\,^{\circ}\mathrm{C}$, resolución 9--12 bits, interfaz 1-Wire, encapsulado estanco. Mide la temperatura del agua. \\
Sensor de corriente & SCT-013 (tipo CT) & No invasivo (pinza), 30/100 A, salida proporcional a la corriente AC. Mide la corriente de la terma para estimar potencia y energía. \\
Sensor de flujo & YF-S201 & 1--30 L/min, $\pm 10\%$, 4,5--24 V, salida de pulsos ($\approx 450$ pulsos/L), conexión 1/2''. Detecta el uso real de agua. \\
Actuador & SSR 25 A (Fotek SSR-25DA o similar) & Optoaislado, control 3--32 V DC, carga 24--380 V AC, conmutación por cruce por cero. Enciende/apaga la terma. \\
Buzzer (opcional) & Módulo buzzer activo & 5 V, alerta sonora local. \\
Acondicionamiento y fuente & Divisor de tensión + filtro RC y fuente AC-DC HLK-PM01 & Adapta las señales a 0--3,3 V; entrega 5 V DC a partir de 220 V AC. \\
Enclosure & Caja ABS con riel DIN (IP54/IP65) & Protección mecánica y eléctrica, separación de la electrónica de baja tensión respecto de la red. \\
\bottomrule
\end{tabularx}
\end{table}

\subsection{Microcontrolador seleccionado}
Se selecciona el \textbf{ESP32-WROOM-32}. La justificación técnica es la siguiente:

\begin{itemize}
  \item \textbf{Conectividad integrada:} Wi-Fi y Bluetooth de forma nativa, lo que elimina módulos externos y reduce costo y complejidad, en línea con las recomendaciones de arquitectura de la literatura~\cite{stojkoska2017,alfuqaha2015}.
  \item \textbf{Capacidad de cómputo:} doble núcleo a 240 MHz y 520 KB de SRAM, suficientes para el procesamiento previo (RMS, filtrado, conteo de pulsos y evaluación de reglas) en el propio nodo~\cite{shi2016}.
  \item \textbf{Periféricos adecuados:} ADC de 12 bits para la señal acondicionada del sensor de corriente, temporizadores por hardware para el conteo de pulsos del sensor de flujo, GPIO para el DS18B20 (1-Wire), la salida de control del SSR y el buzzer.
  \item \textbf{Costo y disponibilidad:} bajo costo y amplia disponibilidad local, requisito explícito del proyecto, en contraste con alternativas más costosas como una computadora de placa única.
\end{itemize}

Alternativas como el ESP8266 (menos periféricos y un solo núcleo) o una Raspberry Pi (mayor costo y consumo) fueron descartadas por no ofrecer la mejor relación capacidad/costo para este prototipo.

\subsection{Diagrama de bloques del sistema propuesto}
La Figura~\ref{fig:bloques} muestra el diagrama de bloques. Los tres sensores entregan sus señales al ESP32, que realiza el acondicionamiento y el procesamiento previo. El ESP32 se comunica bidireccionalmente por Wi-Fi (MQTT) con la nube/dashboard y, a través de éste, con el asistente virtual. Como salida, el ESP32 gobierna el SSR, que conmuta la alimentación de la terma, y opcionalmente un buzzer de alerta.

\begin{figure}[htbp]
\centering
\shorthandoff{<>}
\resizebox{0.98\textwidth}{!}{%
\begin{tikzpicture}[
  font=\footnotesize,
  box/.style={draw, rounded corners=2pt, align=center, inner sep=4pt, minimum height=0.9cm},
  sens/.style={box, fill=primary!8, draw=primary!60, minimum width=3.7cm},
  mcu/.style={box, fill=accent!10, draw=accent!70, minimum width=4.6cm, minimum height=3.6cm, very thick},
  act/.style={box, fill=primary!8, draw=primary!60, minimum width=3.2cm},
  net/.style={box, fill=green!6, draw=green!45!black, minimum width=3.9cm},
  lab/.style={font=\scriptsize\itshape, text=black!65}
]
% Sensores (capa de percepción)
\node[sens] (s1) at (1.9,1.6) {\textbf{Sensor de temperatura}\\ \scriptsize DS18B20 (1-Wire)};
\node[sens] (s2) at (1.9,0)   {\textbf{Sensor de corriente}\\ \scriptsize SCT-013 (no invasivo)};
\node[sens] (s3) at (1.9,-1.6){\textbf{Sensor de flujo}\\ \scriptsize YF-S201 (pulsos)};
% Microcontrolador
\node[mcu] (mcu) at (7.6,0) {\textbf{ESP32}\\[3pt] \scriptsize Acondicionamiento y\\ \scriptsize filtrado de señales\\[2pt] \scriptsize Procesamiento previo\\(RMS, media móvil, pulsos)\\[2pt] \scriptsize Wi-Fi + MQTT};
% Actuadores y carga
\node[act] (ssr) at (13.1,1.7) {\textbf{SSR}\\ \scriptsize (actuador)};
\node[act] (ter) at (17.6,1.7) {\textbf{Terma}\\ \scriptsize 220 V AC};
\node[act] (buz) at (13.1,-1.7) {\textbf{Buzzer}\\ \scriptsize (opcional)};
% Nube y asistente
\node[net] (cld) at (7.6,-3.9) {\textbf{Nube / Broker MQTT}\\ \scriptsize Dashboard};
\node[net] (voz) at (13.1,-3.9) {\textbf{Asistente virtual}\\ \scriptsize (voz / chat)};

% Flechas
\draw[-{Latex}] (s1.east) -- ([yshift=1.0cm]mcu.west);
\draw[-{Latex}] (s2.east) -- (mcu.west);
\draw[-{Latex}] (s3.east) -- ([yshift=-1.0cm]mcu.west);
\draw[-{Latex}] (mcu.east) -- ([yshift=0.9cm]ssr.west) node[midway,above,lab]{control};
\draw[-{Latex}] (ssr.east) -- (ter.west);
\draw[-{Latex}] (mcu.east) -- ([yshift=-0.9cm]buz.west);
\draw[{Latex}-{Latex}] (mcu.south) -- (cld.north) node[midway,right,lab]{Wi-Fi / MQTT};
\draw[{Latex}-{Latex}] (cld.east) -- (voz.west) node[midway,above,lab]{consultas / comandos};
\node[lab, align=center] at (1.9,-2.8) {Capa de percepción};
\node[lab, align=center] at (7.6,2.7) {Capa de red y procesamiento previo};
\node[lab, align=center] at (15.3,3.0) {Capa de actuación};
\end{tikzpicture}%
}
\shorthandon{<>}
\caption{Diagrama de bloques del sistema IoT propuesto para la terma doméstica.}
\label{fig:bloques}
\end{figure}

\subsection{Procesamiento previo y flujo de datos}
El nodo ejecuta el siguiente procesamiento antes de transmitir:
\begin{enumerate}
  \item \textbf{Temperatura:} lectura del DS18B20 y suavizado por media móvil para eliminar ruido y reflejar la inercia térmica del agua.
  \item \textbf{Corriente:} muestreo de la señal acondicionada del SCT-013, cálculo del valor eficaz (RMS) y obtención de la potencia y la energía acumulada.
  \item \textbf{Flujo:} conteo de pulsos por interrupción de hardware, cálculo del caudal instantáneo y del volumen acumulado.
  \item \textbf{Reglas locales:} apagado automático al alcanzar la temperatura objetivo, corte tras un tiempo sin flujo (``encendida sin uso'') y protección ante sobretemperatura.
  \item \textbf{Publicación:} envío de métricas agregadas por MQTT y recepción de comandos de encendido/apagado desde el dashboard o el asistente virtual.
\end{enumerate}

\subsection{Enclosure y consideraciones de seguridad}
La electrónica se aloja en una caja aislante con fijación tipo riel DIN, con separación física entre la etapa de baja tensión (ESP32, sensores y fuente) y la de conmutación de 220 V AC (SSR y borneras). El SSR optoaislado y la fuente AC-DC certificada permiten operar la terma sin contacto eléctrico directo del usuario con la red, atendiendo el requisito de seguridad eléctrica planteado en la problemática.

% ============================================================
\section{Objetivos y Alcances}
\label{sec:objetivos}
% ============================================================

\subsection*{Objetivo general}
Diseñar e implementar un prototipo de sistema IoT inteligente, basado en el microcontrolador ESP32, que permita \textbf{monitorear, controlar de forma remota y gestionar energéticamente una terma eléctrica doméstica}, mediante tres sensores (temperatura, corriente y flujo), un actuador de potencia (SSR), comunicación inalámbrica, un dashboard con procesamiento previo y una interfaz opcional de asistente virtual.

\subsection*{Objetivos específicos}
\begin{enumerate}[label=OE\arabic*.]
  \item Leer de forma simultánea y confiable las variables de temperatura del agua, corriente eléctrica y flujo de agua mediante los tres sensores del sistema.
  \item Acondicionar, filtrar y procesar previamente las señales en el propio nodo (RMS, suavizado, conteo de pulsos) para obtener métricas válidas y reducir el tráfico de red~\cite{shi2016}.
  \item Transmitir la información mediante comunicación inalámbrica (Wi-Fi/MQTT) y visualizarla en un dashboard accesible de forma remota~\cite{alfuqaha2015,elkhozondar2024}.
  \item Accionar de forma remota y automática la carga de la terma mediante un SSR, garantizando el aislamiento eléctrico.
  \item Calcular métricas energéticas (consumo, energía, tiempo de funcionamiento y tiempo de uso) y el estado de calentamiento del agua.
  \item Implementar reglas automáticas de control: apagado al alcanzar la temperatura objetivo y apagado tras un periodo de inactividad sin uso.
  \item Habilitar, de forma opcional y en una segunda etapa, el control mediante comandos de un asistente virtual (``enciende la terma'', ``apaga la terma'', ``¿el agua está caliente?'').
  \item Validar el prototipo experimentalmente, verificando la correspondencia entre las mediciones y el comportamiento real del equipo.
\end{enumerate}

\subsection*{Alcances}
\begin{itemize}
  \item El proyecto entrega un \textbf{prototipo funcional} de terma doméstica inteligente centrado en la terma como caso de uso principal.
  \item El sistema \textbf{mide} temperatura del agua, corriente eléctrica y flujo de agua; \textbf{monitorea} el consumo eléctrico y \textbf{detecta} el uso del agua.
  \item El sistema \textbf{controla} remotamente la carga (encendido/apagado) mediante el SSR y ejecuta \textbf{reglas automáticas} (temperatura objetivo y tiempo sin uso).
  \item El sistema \textbf{transmite} información por comunicación inalámbrica y \textbf{visualiza} el estado en un dashboard, calculando métricas de consumo y funcionamiento.
  \item Se contempla una \textbf{interfaz opcional} de asistente virtual mediante comandos de voz o chat.
  \item El alcance incluye el \textbf{procesamiento previo} de los datos en el nodo y el \textbf{enclosure} del prototipo.
\end{itemize}

\subsection*{Limitaciones}
\begin{itemize}
  \item El prototipo se valida en un entorno controlado (banco de pruebas / vivienda piloto), no como producto comercial listo para producción.
  \item El asistente virtual es una funcionalidad \textbf{opcional y de segunda prioridad}: se avanzará de forma paralela, pero sus entradas sólo se considerarán fiables una vez que el sensado y la actuación estén operativos.
  \item La interfaz con el asistente virtual depende de servicios de reconocimiento de voz y de conectividad a Internet; su disponibilidad no está garantizada en todos los entornos.
  \item La medición de corriente con CT es de tipo estimación (no de laboratorio); no se busca una precisión de instrumentación certificada.
  \item El control es de tipo encendido/apagado (on-off) sobre la carga; no se implementa regulación proporcional de potencia (dimmer/control de fase), por simplicidad y seguridad.
  \item La solución se restringe a una terma monofásica de 220 V AC; no se contemplan instalaciones trifásicas ni otras cargas.
  \item No se abordan en esta etapa aspectos avanzados de ciberseguridad (por ejemplo, gestión de certificados), aunque se usan credenciales y red segura.
\end{itemize}

% ============================================================
\section*{Conclusiones}
\addcontentsline{toc}{section}{Conclusiones}
% ============================================================
El perfil de proyecto define una solución IoT viable, económica y técnicamente fundamentada para el monitoreo, control remoto y gestión energética de una terma doméstica. La propuesta articula tres sensores, un actuador de potencia seguro, un microcontrolador ESP32 con comunicación inalámbrica, procesamiento previo en el nodo y una interfaz con dashboard y asistente virtual. El estado del arte confirma la pertinencia del enfoque y evidencia el vacío que el proyecto busca cubrir. La prioridad de implementación será asegurar la lectura confiable de los tres sensores, la transmisión inalámbrica y el accionamiento remoto y automático del SSR; el asistente virtual se desarrollará de forma paralela y se habilitará una vez estabilizadas las entradas de datos.

% ============================================================
% REFERENCIAS
% ============================================================
\clearpage
\section*{Referencias}
\addcontentsline{toc}{section}{Referencias}
\begin{thebibliography}{99}\small

\bibitem{alfuqaha2015}
Al-Fuqaha, A., Guizani, M., Mohammadi, M., Aledhari, M., \& Ayyash, M. (2015).
Internet of Things: A survey on enabling technologies, protocols, and applications.
\emph{IEEE Communications Surveys \& Tutorials}, 17(4), 2347--2376.

\bibitem{shi2016}
Shi, W., Cao, J., Zhang, Q., Li, Y., \& Xu, L. (2016).
Edge computing: Vision and challenges.
\emph{IEEE Internet of Things Journal}, 3(5), 637--646.

\bibitem{zoha2012}
Zoha, A., Gluhak, A., Imran, M. A., \& Rajasegarar, S. (2012).
Non-intrusive load monitoring approaches for disaggregated energy sensing: A survey.
\emph{Sensors}, 12(12), 16838--16866.

\bibitem{alaa2017}
Alaa, M., Zaidan, A. A., Zaidan, B. B., Talal, M., \& Kiah, M. L. M. (2017).
A review of smart home applications based on Internet of Things.
\emph{Journal of Network and Computer Applications}, 97, 48--65.

\bibitem{stojkoska2017}
Stojkoska, B. L. R., \& Trivodaliev, K. V. (2017).
A review of Internet of Things for smart home: Challenges and solutions.
\emph{Journal of Cleaner Production}, 140, 1454--1464.

\bibitem{bhati2017}
Bhati, A., Hansen, M., \& Chan, C. M. (2017).
Energy conservation through smart homes in a smart city: A lesson for Singapore households.
\emph{Energy Policy}, 104, 230--239.

\bibitem{ford2017}
Ford, R., Pritoni, M., Sanguinetti, A., \& Karlin, B. (2017).
Categories and functionality of smart home technology for energy management.
\emph{Building and Environment}, 123, 543--554.

\bibitem{marikyan2019}
Marikyan, D., Papagiannidis, S., \& Alamanos, E. (2019).
A systematic review of the smart home literature: A user perspective.
\emph{Technological Forecasting and Social Change}, 138, 139--154.

\bibitem{zipperer2013}
Zipperer, A., Aloise-Young, P. A., Suryanarayanan, S., Roche, R., Earle, L., Christensen, D., \& Zimmerle, D. (2013).
Electric energy management in the smart home: Perspectives on enabling technologies and consumer behavior.
\emph{Proceedings of the IEEE}, 101(11), 2397--2408.

\bibitem{elkhozondar2024}
El-Khozondar, H. J., Mtair, S. Y., Qoffa, K. O., Qasem, O. I., et al. (2024).
A smart energy monitoring system using ESP32 microcontroller.
\emph{e-Prime -- Advances in Electrical Engineering, Electronics and Energy}.

\bibitem{ds18b20}
Maxim Integrated. (2019).
\emph{DS18B20 Programmable Resolution 1-Wire Digital Thermometer} (Datasheet).

\bibitem{ssr}
Fotek Controls.
\emph{SSR-25DA Solid State Relay} (Datasheet).

\bibitem{yfs201201}
Seeed Studio.
\emph{YF-S201 Water Flow Sensor} (Datasheet).

\end{thebibliography}

\end{document}
